% HD-F1-T20. GPT-6 Astra Pro; /root/transport; 22 November 2026 UTC.
\section[HD-F1-T20: exclusion of symplectic multiplicity seven]{HD-F1-T20: the original tensor inertia excludes symplectic-type multiplicity seven}
\label{sec:f1-symplectic-seventh-multiplicity}

Retain the actual normal subgroup, faithful scalar centre,
nondegenerate commutator pairing and original coefficient
representation of T17. The multiplicity of its irreducible
constituent is at least eight. In particular all three
remaining multiplicity-seven dimensions of T17 are excluded.
We use the already proved actual evaluation isomorphism
\[
 \epsilon:W\otimes M\xrightarrow{\sim}V_{\mathbf C},
 \qquad d=\dim W=2^a,\quad k=\dim M,
 \quad C=C_G(R),
 \tag{F1T20.1}
\]
and, for this proof, \(k=7\) and
\(m=7d\in\{28,56,112\}\). Every nonidentity
element of the remaining T2 coefficient image has moving
rank at least five. The original tuple still has
\(b\ge12\), at least three entries outside \(C\),
product one, and total moving rank \(2m+6\).

\paragraph{1. Actual inertia intertwiners and their scalar choices.}
For each original inertia \(g\), conjugation by \(g\)
fixes the faithful central character of \(R\).
The averaged-Hom uniqueness argument in T17 therefore
gives an invertible intertwiner \(A_g:W\to W\) satisfying
\[
 A_g\rho_W(r)A_g^{-1}=\rho_W(grg^{-1})
 \quad(r\in R).
 \tag{F1T20.2}
\]
The operators \(\epsilon^{-1}g\epsilon\) and
\(A_g\otimes1\) induce the same conjugation on
\(\rho_W(R)\otimes1\). Their quotient centralizes this
action, so the actual centralizer computation in T17 gives
a unique \(T_g\in\operatorname{GL}(M)\) with
\[
       \epsilon^{-1}g\epsilon=A_g\otimes T_g.
 \tag{F1T20.3}
\]
The permitted scalar change is exactly
\((A_g,T_g)\mapsto(zA_g,z^{-1}T_g)\).
For \(g\in C\) we choose \(A_g=1\), leaving the
actual centralizer action \(T_g\) of T17.

These factors can be chosen semisimple without any
replacement of the original matrix. If \(g^N=1\),
then \(A_g^N\) centralizes \(R\) and is a scalar by
Schur's lemma on \(W\). Rescale \(A_g\) so that
\(A_g^N=1\); equation (F1T20.3) then gives
\(T_g^N=1\). All subsequent scalar changes preserve
semisimplicity. For \(g\notin C\), \(A_g\) is not
scalar. The exact commutator and perfect-pairing argument
of T15, now on the actual irreducible \(R\)-space \(W\),
applies to every eigenspace of this intertwiner. It gives
\[
 \dim\ker(A_g-\lambda)\le\frac{3d}{4}
 \quad\hbox{for every }\lambda\in\mathbf C.
 \tag{F1T20.4}
\]
Explicitly, a noncentral commutator has every eigenspace
of dimension \(d/2\), while
\([A_g,e]-1=((A_g-\lambda)e-e(A_g-\lambda))A_g^{-1}e^{-1}\)
has rank at most twice that of \(A_g-\lambda\).
When all commutators are central, the perfect pairing
gives \(A_g=z\rho_W(u)\) for a noncentral \(u\in R\),
whose two eigenspaces both have dimension \(d/2\).
Thus the same bound also covers that case.

\paragraph{2. The exact cost of a nonscalar multiplicity factor.}
Let \(a_1\ge a_2\ge\cdots\) be the eigenspace
dimensions of \(A_g\), and pair them with the dimensions
of the inverse eigenspaces of \(T_g\). Unmatched
eigenvalues contribute zero. This is the exact fixed-space
decomposition of \(A_g\otimes T_g\). If \(T_g\)
is nonscalar, its at least two positive eigenspace
dimensions sum to seven. Consequently
\[
 \begin{split}
 \dim V^g&\le6a_1+a_2
       \le d+5a_1\le\frac{19d}{4},\\
 \operatorname{rank}(g-1)&\ge\frac{9d}{4}.
 \end{split}
 \tag{F1T20.5}
\]
Equality forces exactly two eigenspaces of \(A_g\),
of dimensions \(3d/4,d/4\), and two eigenspaces of
\(T_g\), of dimensions \(6,1\), paired in that
order by inverse eigenvalues. These equality assertions
follow from equality in each displayed inequality:
\(a_1=3d/4\), \(a_2=d-a_1\), and all seven
multiplicity dimensions are assigned as six to the
strictly largest \(a_1\) and one to \(a_2\).
If \(T_g\) is scalar, the moving rank of \(g\)
is a multiple of seven and is at least \(7d/4\).

\paragraph{3. The dimensions 28 and 112.}
For \(m=28\), the moving rank inside \(C\) is a
multiple of \(d=4\). The original no-rank-at-most-four
gate therefore makes it at least eight. Outside \(C\)
the rank is at least seven. Twelve original branches
cost at least \(84>2\cdot28+6\), a contradiction.

For \(m=112\), write \(n\) for the number of
outside-centralizer entries. The inside and outside
lower bounds give
\[
       230\ge16b+12n,\qquad b\ge12,\quad n\ge3.
 \tag{F1T20.6}
\]
Thus \(b=12\) and \(n=3\). The lower cost is
\(228\), and an increase in any inside rank costs
at least sixteen. All nine inside ranks are therefore
exactly sixteen; the three outside ranks sum to 86,
and each lies between 28 and 30. A nonscalar \(T_g\)
would cost at least 36 by (F1T20.5), so all three
are scalar. Their ranks are multiples of seven and
are therefore all 28. Their sum is 84, contradicting 86.

\paragraph{4. The dimension 56 and its actual product-one factors.}
For \(m=56\) the same argument gives
\[
       118\ge8b+6n,\qquad b\ge12,\quad n\ge3,
 \tag{F1T20.7}
\]
so again \(b=12,n=3\). All nine inside ranks are
exactly eight, and the three outside ranks sum to 46,
with each between 14 and 18. A scalar \(T_g\) makes
its rank a multiple of seven, hence 14. A nonscalar
\(T_g\) costs at least 18 by (F1T20.5). The only
possible outside pattern is therefore \(14,14,18\).
The nine inside matrices have \(A_g=1\) and
\(\operatorname{rank}(T_g-1)=1\).

For either rank-14 outside matrix, write \(T_g=t1\)
and replace its factors by \((tA_g,1)\). Then
\(\operatorname{rank}(A_g-1)=2\). For the rank-18
matrix, equality in (F1T20.5) gives eigenspace
dimensions \((6,2)\) on \(W\) and \((6,1)\)
on \(M\), with inverse matching. Let \(a\) be
the eigenvalue of \(A_g\) on its six-dimensional
eigenspace. Replace its factors by
\((a^{-1}A_g,aT_g)\). The eigenvalue one then
has multiplicity six in both factors, so their
moving ranks are respectively two and one.

These scalar changes preserve every original tensor
matrix, including their order in the original tuple.
The product-one relation gives
\[
 (A_1\cdots A_{12})\otimes(T_1\cdots T_{12})=1.
 \tag{F1T20.8}
\]
An eigenvector of the second invertible factor shows
that the first is a scalar, say \(z1_W\); then the
second is \(z^{-1}1_M\). Only three \(A_j\) differ
from one and each has moving rank two. The exact
telescoping rank inequality yields
\[
 \operatorname{rank}(z1_W-1)
 \le\sum_j\operatorname{rank}(A_j-1)=6<8.
 \tag{F1T20.9}
\]
Hence \(z=1\), and the actual chosen tuple on
\(M\) has product one. Its total moving rank is
\(9+1=10\).

\paragraph{5. The rank inequality on this actual tuple.}
The tuple \(T_1,\ldots,T_{12}\) is irreducible on
the seven-dimensional \(M\). A common invariant
subspace \(U\subset M\) would give, under
\(\epsilon\), a common invariant subspace
\(W\otimes U\subset V_{\mathbf C}\) for all the
original \(g_j\); they generate the original
absolutely irreducible \(G\). Thus \(U=0\) or
\(M\). Both the common invariant space of the
tuple and that of its dual are zero, since \(\dim M>1\).

Here is the full cochain proof of the needed Scott
inequality. Set \(P_0=1\), \(P_j=T_1\cdots T_j\),
and form the actual linear maps
\[
 \begin{aligned}
 F:M&\longrightarrow\bigoplus_j\operatorname{Im}(T_j-1),\\
 x&\longmapsto((T_j-1)x)_j,\\
 D:\bigoplus_j\operatorname{Im}(T_j-1)&\longrightarrow M,\\
 (v_j)_j&\longmapsto\sum_jP_{j-1}v_j.
 \end{aligned}
 \tag{F1T20.10}
\]
The kernel of \(F\) is the zero common invariant
space. The annihilator of \(\operatorname{Im}D\)
consists exactly of the common invariant covectors:
\(\varphi P_{j-1}(T_j-1)=0\) implies successively
\(\varphi P_j=\varphi P_{j-1}=\varphi\).
It is zero, so \(D\) is surjective. Finally,
\(DF=P_{12}-1=0\) by the proved product-one
identity. Consequently
\[
 \sum_j\operatorname{rank}(T_j-1)
   =\dim\ker D+\dim M\ge2\dim M=14,
 \tag{F1T20.11}
\]
contradicting ten. This excludes the final dimension.
Together with T17, multiplicities one through seven
are now all excluded at the original source.

\paragraph{Original source and geometric framework.}
The maps used here are the original isotypic evaluation,
the explicit inertia intertwiners and their inverse scalar
changes, the fixed-space tensor decomposition, and the
cochain maps (F1T20.10). The tuple on \(M\) is used
only in this proved matrix contradiction; no Jacobian
factor is replaced by a putative multiplicity source.
The connected cover, coefficient sheaf, rational primitive
projector, actual sixfold corner and restricted theta
form are those of T17. Their Poincare biextension,
parabolic cup pairing, Hodge filtration, period map,
Grothendieck group and Adams realizations carry their
identity maps with multiplier one. The Scott rank
inequality is classical; the complete required proof
and its exact scalar-to-product-one map are supplied
above. The new conclusion is the exclusion of these
three actual source cases, with no extra detector
component counted.
